\chapter{Usage of Chatbots}
\label{chap:tools_ai_usage}

The integration of chatbots powered by advanced language models has become increasingly useful in project-based work. Such tools can provide rapid, context-aware assistance in technical and writing-related tasks, help identify possible issues, and support the preparation of structured documentation. Their ability to respond to natural-language queries and generate relevant suggestions can improve efficiency during development and reporting, provided that their output is reviewed critically and not used without verification.

In the context of this project, chatbot-based and AI-assisted tools were used in a limited and supportive role. In particular, ChatGPT and Gemini were used for several tasks during the development of the code and the preparation of the report. Grammarly was also used as an auxiliary language-support tool during the writing process.

The main areas of use were the following:
\begin{itemize}
    \item \textbf{Code review:} Chatbot support was used to review parts of the MATLAB code, identify possible weaknesses, and suggest improvements in structure, clarity, or consistency.

    \item \textbf{Debugging support:} The tools were used to discuss possible causes of errors, unexpected behavior, or implementation problems and to generate suggestions for resolving them.

    \item \textbf{Code comments:} Assistance was used in formulating or improving code comments so that the implementation became easier to read and understand.

    \item \textbf{Report writing in \LaTeX:} Chatbot support was used to help formulate technical text and convert report content into well-structured \LaTeX{} code during the preparation of the written report.

    \item \textbf{Language support:} Grammarly was used in a limited way to support wording, grammar, and language refinement in the final report.
\end{itemize}

The use of these tools was restricted to supportive assistance only. No chatbot-generated code was adopted directly without modification. Instead, suggestions were reviewed, adapted, and rewritten before being incorporated into the project. Likewise, no generated report text was inserted unchanged into the final document; all written content was manually revised and edited before inclusion in the final submission. All parts of the project, including the implementation, results, code comments, and report text, were manually checked before inclusion in the final submission.

Overall, the use of chatbots and AI-assisted language tools improved the efficiency of selected parts of the workflow, especially in code review, debugging discussion, commenting, and report preparation. However, the final technical decisions, implementation, verification, interpretation of results, and wording of the report remained entirely under manual responsibility.